\documentclass[a4paper,11pt]{article}
\usepackage{booktabs,graphicx,amsmath,amssymb,hyperref,url}
\title{Supplementary Material for CRC-GAD}
\author{}
\date{}
\begin{document}
\maketitle

\section{Per-seed main results}
Table~\ref{tab:per_seed} lists every seed of the primary injected-benchmark protocol.

\begin{table}[ht]
\centering
\caption{Per-seed results from \texttt{results/main\_results.csv}.}
\label{tab:per_seed}
\scriptsize
\begin{tabular}{llcccc}
\toprule
Dataset & Seed & Raw AUC & Conf.\ AUC & FPR@0.05 & TPR@0.05 \\
\midrule
acm & 0 & 0.666 & 0.672 & 0.045 & 0.068 \\
acm & 1 & 0.627 & 0.628 & 0.055 & 0.084 \\
acm & 2 & 0.675 & 0.685 & 0.051 & 0.095 \\
acm & 3 & 0.670 & 0.680 & 0.052 & 0.059 \\
acm & 4 & 0.745 & 0.740 & 0.050 & 0.134 \\
blogcatalog & 0 & 0.637 & 0.658 & 0.062 & 0.030 \\
blogcatalog & 1 & 0.633 & 0.659 & 0.056 & 0.013 \\
blogcatalog & 2 & 0.649 & 0.654 & 0.038 & 0.062 \\
blogcatalog & 3 & 0.672 & 0.670 & 0.053 & 0.049 \\
blogcatalog & 4 & 0.633 & 0.640 & 0.042 & 0.000 \\
citeseer & 0 & 0.802 & 0.785 & 0.031 & 0.512 \\
citeseer & 1 & 0.796 & 0.741 & 0.019 & 0.385 \\
citeseer & 2 & 0.798 & 0.833 & 0.032 & 0.605 \\
citeseer & 3 & 0.824 & 0.800 & 0.031 & 0.452 \\
citeseer & 4 & 0.819 & 0.785 & 0.008 & 0.400 \\
cora & 0 & 0.786 & 0.740 & 0.026 & 0.194 \\
cora & 1 & 0.771 & 0.821 & 0.032 & 0.484 \\
cora & 2 & 0.783 & 0.853 & 0.057 & 0.618 \\
cora & 3 & 0.780 & 0.789 & 0.052 & 0.500 \\
cora & 4 & 0.761 & 0.636 & 0.031 & 0.103 \\
flickr & 0 & 0.606 & 0.619 & 0.048 & 0.000 \\
flickr & 1 & 0.662 & 0.655 & 0.051 & 0.000 \\
flickr & 2 & 0.630 & 0.618 & 0.048 & 0.000 \\
flickr & 3 & 0.632 & 0.654 & 0.060 & 0.000 \\
flickr & 4 & 0.646 & 0.648 & 0.061 & 0.000 \\
pubmed & 0 & 0.681 & 0.685 & 0.037 & 0.426 \\
pubmed & 1 & 0.821 & 0.815 & 0.047 & 0.177 \\
pubmed & 2 & 0.729 & 0.724 & 0.035 & 0.517 \\
pubmed & 3 & 0.754 & 0.718 & 0.036 & 0.166 \\
pubmed & 4 & 0.838 & 0.835 & 0.041 & 0.137 \\
\bottomrule
\end{tabular}

\end{table}

\section{CoLA reference vs.\ pinned implementation}
Section~4 of the main paper states the CoLA reference formulation (RWR subgraphs, bilinear discriminator).
Our pinned code (\texttt{src/crc\_gad/scorer.py}) uses a simplified NumPy GCN with neighbor-mean positive context and shuffled negative context.
All primary numbers use the simplified scorer; the reference formulation is cited for attribution, not as a claim that we re-ran the original PyTorch CoLA binary.

\section{Guarantee taxonomy (restated)}
\begin{itemize}
  \item Proved: marginal $\Pr(p_i\leq\alpha)\leq\alpha$ for a uniformly random test node under a random partition.
  \item Not proved as a theorem: $\mathbb{E}[\widehat{\mathrm{FPR}}]\leq\alpha$ for normals without a ranking/dominance bridge.
  \item Not claimed: node-conditional or community-conditional FPR control.
\end{itemize}

\section{PyTorch DOMINANT backbone}
We implement Ding et al.'s DOMINANT in PyTorch (\texttt{src/crc\_gad/dominant\_torch.py}) following the authors' public layout: shared 2-layer GCN encoder, GCN attribute decoder, structure decoder $A{=}ZZ^\top$, and score $\alpha\cdot e_{\mathrm{attr}}+(1-\alpha)\cdot e_{\mathrm{struct}}$.
Training uses full-batch Adam for 100 epochs ($\mathrm{hid}{=}64$, $\mathrm{lr}{=}5\cdot 10^{-3}$, dropout $0.3$, $\alpha{=}0.8$).
CRC-GAD wraps the frozen scores under the same label-free $\mathcal{C}/\mathcal{T}$ split as other backbones.
Results: \texttt{results/dominant\_pytorch.csv}; regenerate via \texttt{scripts/run\_dominant\_pytorch.py}.

\section{Extended notes on contamination and bridges}
\label{sec:supp_contamination}
When calibration contains a fraction $\pi_{\mathcal{C}}$ of anomalies and anomalous scores stochastically dominate normals, normal test p-values tend to inflate (conservative FPR).
If the scorer inverts the ranking, contamination can be anti-conservative.
The main paper's Theorem on random-test-node validity does not by itself imply $\mathbb{E}[\widehat{\mathrm{FPR}}]\leq\alpha$; that step needs the dominance bridge stated in the main text.
Degree-stratified FPR (main paper) shows that marginal validity can coexist with subgroup FPR far from $\alpha$.

\section{Extended studies}
Additional CSV: \texttt{results/extended\_studies.csv} (multi-backbone, multi-dataset heuristics, degree FPR, expanded contamination, organic proxy), regenerated by \texttt{scripts/run\_extended\_studies.py}.

\end{document}
